%=====================================================================
\section{Late rules and unique link keys}\label{sec:v75latekeys}
% numbered from 1880; the shared counter is restored at the end
\newcounter{savedproplatekeys}\setcounter{savedproplatekeys}{\value{proposition}}
Two classes of declared constraint were reported but not always met. A \emph{late rule} is a check rule that reads an aggregate column, on its own table or through a link (a card issued on or after its account's opening date, when the opening date is the account's first transaction date). It can be evaluated only after every table is drawn and the aggregates computed. A \emph{unique link key} is a unique key made of one link and columns fixed before the links (one owner per account: \texttt{[account\_id, type]}, with the type set from the disposition's stratum). When nearly every parent already holds the key, the link model's proposals and a uniform fallback rarely reach the few free parents, so rows were left in conflict. The code is \texttt{Synthesizer.\_enforce\_late} and \texttt{Synthesizer.\_free\_parents}; the tests are in \texttt{tests/unit/test\_late\_rules\_unique\_links.py}.

\setcounter{proposition}{1879}
\begin{proposition}[Free-parent draws keep the link model on the valid set]\label{prop:freeparents}
Let a group of $m$ rows share the values of the other columns of every unique key that holds link $L$, and let $F$ be the set of parents whose key is not held by another row of the group. Row $i$ of segment $k_i$ draws parent $p$ with probability proportional to $\pi_i(p)=\phi_{k_i,z_p}\,w_p/m_{z_p}$, the link model's probability (segment link law $\phi$, degree weight $w_p$, $m_s$ parents in segment $s$). The rows of one segment draw distinct parents from $F$ by successive sampling proportional to $\pi$ (Gumbel top-$k$), and the parents taken are removed from $F$ before the next segment's rows. (a) When $|F|\ge m$ and some parent of $F$ has positive weight for every segment (else the draw is uniform on $F$), every row of the group is placed and the group's key stays unique. (b) Each placement has the law $\pi_i$ restricted to the parents still free, i.e.\ the link model conditioned on the key constraint for that row given the rows already placed. (c) Rows that already satisfied every key and rule are not moved.
\end{proposition}
\begin{proof}
(a) Each segment's block takes $\min(\text{rows},|\{p\in F:\pi(p)>0\}|)$ distinct parents of $F$ and removes them; with positive weight on $F$ for each segment (or the uniform fallback) and $|F|\ge m$, every row receives a distinct free parent, so no two rows of the group hold the same key and no row takes a key held outside the group. (b) Adding independent Gumbel noise to $\log\pi$ and taking the top $r$ is equivalent to drawing $r$ items one at a time without replacement with probability proportional to $\pi$ among those left (Vieira 2014; Kool et al.\ 2019), so each draw is $\pi_i$ restricted to the parents still free. (c) The repair runs only on rows marked invalid, and the move is accepted by the same check as every other move of the unique-key sampler (key rules, mandatory participation, keys). The MH sweeps that follow keep the restricted law invariant (Section on unique keys).
\end{proof}
\noindent The draw is a repair, as was the uniform fallback it replaces for links in a unique key: it reaches a valid state with the link model's weights where the uniform fallback ignored them. Setting \texttt{uniqueness.free\_parent\_repair: off} restores the uniform fallback. Test \texttt{test\_unique\_link\_key\_exact}.

\setcounter{proposition}{1880}
\begin{proposition}[Late rules by redraw after the aggregates]\label{prop:laterules}
Let $R$ be the late rules of table $t$, $Q(x\mid z)$ the generator's law of a row's modelled attributes given its segment (and conditioning values), and $V_i$ the set of attribute values for which row $i$ satisfies $R$, keeps every other rule of $t$ it satisfied, keeps the unique keys over attributes, and is not a near copy. Rows that break $R$ draw $x\sim Q(\cdot\mid z_i)$ until $x\in V_i$, at most \texttt{rules.max\_rejection\_rounds} times; the aggregates, lookups and derives are then recomputed. (a) An accepted row has the law $Q(\cdot\mid z_i,\,x\in V_i)$. (b) A row is changed only when it broke a late rule and the new values satisfy all of them, so no rule that held on a row is broken by the redraw. (c) When the redrawn table does not feed an aggregate, $V_i$ does not change during the redraws and one pass suffices; otherwise the pass is repeated on the recomputed aggregates, at most three times, and the remaining violations are reported.
\end{proposition}
\begin{proof}
(a) Rejection sampling: draws are independent from $Q(\cdot\mid z_i)$ and accepted exactly on $V_i$, so the accepted value has the restricted law (the argument of the rule redraws in \texttt{\_finish}). (b) Acceptance requires every late rule and every rule that held before. The redraw is made only when no rule of another table reads the redrawn columns, and the table is not one whose other values are computed from these attributes (versions, sequences, anchored links, cumulative rules, scenario maps); the final violation counts are recomputed on all tables after the aggregates. (c) $V_i$ depends on the aggregates only through the columns the rules read; if $t$ is not aggregated, these do not change while $t$ is redrawn.
\end{proof}
\noindent Privacy: the redrawn values are new draws from the fitted model, as for any rule redraw, and are subject to the near-copy guard where it applies; no source row is read at sampling time. \texttt{rules.late: report} keeps the earlier behaviour (late rules counted, not enforced). Test \texttt{test\_late\_rule\_enforced}.

\setcounter{proposition}{\value{savedproplatekeys}}
